\section{Methodology}
\label{sec:method}

\subsection{System Architecture Overview}

Figure~\ref{fig:architecture} illustrates the three-phase closed-loop architecture. Each review cycle proceeds through synthesis → planning → execution → validation.

\begin{figure}[t]
\centering
\begin{verbatim}
   ┌─────────────┐
   │  Reviews    │ REVIEW-ALICE.md, REVIEW-BOB.md, ...
   └──────┬──────┘
          ↓
   ┌─────────────┐
   │  Synthesis  │ SYNTHESIS.md (P1/P2/P3 items)
   └──────┬──────┘
          ↓
   ┌─────────────┐
   │  Planning   │ REVISION-PLAN.md (file, line, old, new)
   └──────┬──────┘
          ↓
   ┌─────────────┐
   │  Execution  │ Apply edits to LaTeX source
   └──────┬──────┘
          ↓
   ┌─────────────┐
   │  Validation │ Compile LaTeX, re-review, repeat
   └─────────────┘
\end{verbatim}
\caption{Three-phase closed-loop architecture from reviews to revisions.}
\label{fig:architecture}
\end{figure}

\subsection{Phase 1: Review Synthesis}

Review synthesis follows methods from prior work~\cite{dellalibera2026synthesis}:

\textbf{Inputs}: 5-7 review files (REVIEW-EXPERT-NAME.md) with structured feedback (strengths, weaknesses, questions, score).

\textbf{Process}:
\begin{enumerate}
\item Extract all issues from weaknesses sections.
\item Deduplicate issues using semantic similarity (cosine similarity on sentence embeddings, threshold 0.85).
\item Classify priority based on reviewer consensus:
  \begin{itemize}
  \item P1: 3+ reviewers flag issue OR issue threatens acceptance
  \item P2: 2+ reviewers flag issue
  \item P3: 1 reviewer flags issue
  \end{itemize}
\item Generate SYNTHESIS.md with prioritized issue list.
\end{enumerate}

\textbf{Outputs}: SYNTHESIS.md with P1/P2/P3 sections listing issues in order of priority.

\subsection{Phase 2: Revision Planning}

Revision planning translates abstract feedback into concrete edit instructions.

\textbf{Inputs}:
\begin{itemize}
\item SYNTHESIS.md with prioritized issues
\item Paper LaTeX source (main.tex, sections/*.tex)
\item Prior REVISION-PLAN.md (if round ≥ 2) to avoid redundant edits
\end{itemize}

\textbf{Process}: For each issue in SYNTHESIS.md (P1 first, then P2, then P3):
\begin{enumerate}
\item \textbf{Localize}: Identify the file and line range where the issue occurs. Use grep to search for relevant keywords (e.g., issue "methodology section lacks detail" → search for "methodology" in sections/*.tex).

\item \textbf{Propose edit}: Generate old text (exact text to replace) and new text (replacement). Ensure old text is unique within the file to avoid ambiguous matches.

\item \textbf{Validate}: Check that old text exists in the file at the specified location. If not found, mark issue as "requires human judgment" (cannot localize automatically).

\item \textbf{Format}: Write edit as structured entry in REVISION-PLAN.md:
\begin{verbatim}
## P1-3: Methodology lacks statistical power analysis

**File**: sections/03-methodology.tex
**Lines**: 127-132
**Action**: Insert

**Old text**: (none - insertion)

**New text**:
We conducted a power analysis using G*Power 3.1 \cite{faul2007gpower}.
With $\alpha = 0.05$ and desired power $1-\beta = 0.80$, the minimum
sample size is $n = 47$ for detecting a medium effect size ($d = 0.5$).

**Rationale**: Reviewers Alice, Bob, and Carol all noted the lack of
statistical justification for sample size. Adding power analysis
addresses this P1 concern.
\end{verbatim}
\end{enumerate}

\textbf{Outputs}: REVISION-PLAN.md with structured edit entries for all addressable issues. Issues requiring human judgment (e.g., "restructure paper", "expand related work with 3+ new citations") are marked but not assigned concrete edits.

\subsection{Phase 3: Edit Execution}

Edit execution parses REVISION-PLAN.md and applies edits to LaTeX files.

\textbf{Inputs}:
\begin{itemize}
\item REVISION-PLAN.md with structured edit entries
\item Paper LaTeX source (writable)
\end{itemize}

\textbf{Process}:
\begin{enumerate}
\item Parse REVISION-PLAN.md to extract edit entries (file, action, old text, new text).

\item For each edit entry in priority order (P1 → P2 → P3):
  \begin{enumerate}
  \item Read target file.
  \item Verify old text exists (exact match). If not found, skip edit and log warning.
  \item Apply edit using string replacement:
    \begin{itemize}
    \item \textbf{Replace}: Substitute old text with new text.
    \item \textbf{Insert}: Add new text at specified line.
    \item \textbf{Delete}: Remove old text.
    \end{itemize}
  \item Write modified file.
  \item Mark edit as completed in \_panel.yaml state file.
  \end{enumerate}

\item After all edits applied, run \texttt{make pdf} to compile LaTeX.

\item If compilation fails:
  \begin{itemize}
  \item Identify the failing edit (last edit before error).
  \item Rollback that edit (restore old text).
  \item Mark edit as "requires human judgment".
  \item Recompile to verify rollback succeeded.
  \end{itemize}

\item If compilation succeeds, commit changes to git with message listing completed P1 items.
\end{enumerate}

\textbf{Outputs}:
\begin{itemize}
\item Modified LaTeX files with edits applied
\item Compiled PDF (if successful)
\item \_panel.yaml updated with completion status per revision item
\end{itemize}

\subsection{Phase 4: Validation and Iteration}

After edits are applied, the paper enters the recheck stage for re-review.

\textbf{Process}:
\begin{enumerate}
\item Compile paper to PDF.
\item Generate round 2 reviews (ROUND2-REVIEW-EXPERT-NAME.md) from the same reviewers.
\item Synthesize round 2 feedback into new SYNTHESIS.md.
\item Check convergence criteria:
  \begin{itemize}
  \item Average score ≥ 2.5/4 (weak accept or better)
  \item No individual score < 2/4 (no rejects)
  \item All P1 items from prior round addressed
  \end{itemize}
\item If converged: paper advances to ready stage (panel tier review).
\item If not converged: generate new REVISION-PLAN.md and repeat phases 2-4.
\end{enumerate}

Iteration continues until convergence or a maximum of 4 rounds (empirically, 85\% of papers converge within 2 rounds~\cite{dellalibera2026dynamics}).

\subsection{Edit Types and Complexity}

We classify edits by complexity:

\textbf{Simple edits} (fully automatable):
\begin{itemize}
\item Phrase replacement: "replace X with Y"
\item Insertion: "add text at line N"
\item Deletion: "remove text at line N"
\item Typo fixes: "change 'recieve' to 'receive'"
\end{itemize}

\textbf{Complex edits} (require human judgment):
\begin{itemize}
\item Structural changes: "move section 4 before section 3"
\item Content expansion: "add 3 paragraphs explaining X"
\item Figure creation: "add figure showing Y"
\item Major rewrites: "rewrite introduction for clarity"
\end{itemize}

The system attempts all simple edits automatically. Complex edits are flagged in REVISION-PLAN.md with \texttt{[HUMAN REQUIRED]} tags, signaling that author intervention is needed.

\subsection{Implementation}

The system is implemented as part of the panel plugin for Claude Code:

\begin{itemize}
\item \textbf{Synthesis}: \texttt{panel:paper} command, synthesis stage
\item \textbf{Planning}: \texttt{panel:paper} command, revision stage (plan generation)
\item \textbf{Execution}: \texttt{panel:paper} command, revision stage (edit application)
\item \textbf{Validation}: \texttt{panel:paper} command, recheck stage
\end{itemize}

Edits use Claude Code's Edit tool, which performs exact string matching and diff validation, ensuring edits are applied precisely as specified or fail with clear error messages.